% Options for packages loaded elsewhere
\PassOptionsToPackage{unicode}{hyperref}
\PassOptionsToPackage{hyphens}{url}
\documentclass[
  ignorenonframetext,
]{beamer}
\newif\ifbibliography
\usepackage{pgfpages}
\setbeamertemplate{caption}[numbered]
\setbeamertemplate{caption label separator}{: }
\setbeamercolor{caption name}{fg=normal text.fg}
\beamertemplatenavigationsymbolsempty
% remove section numbering
\setbeamertemplate{part page}{
  \centering
  \begin{beamercolorbox}[sep=16pt,center]{part title}
    \usebeamerfont{part title}\insertpart\par
  \end{beamercolorbox}
}
\setbeamertemplate{section page}{
  \centering
  \begin{beamercolorbox}[sep=12pt,center]{section title}
    \usebeamerfont{section title}\insertsection\par
  \end{beamercolorbox}
}
\setbeamertemplate{subsection page}{
  \centering
  \begin{beamercolorbox}[sep=8pt,center]{subsection title}
    \usebeamerfont{subsection title}\insertsubsection\par
  \end{beamercolorbox}
}
% Prevent slide breaks in the middle of a paragraph
\widowpenalties 1 10000
\raggedbottom
\AtBeginPart{
  \frame{\partpage}
}
\AtBeginSection{
  \ifbibliography
  \else
    \frame{\sectionpage}
  \fi
}
\AtBeginSubsection{
  \frame{\subsectionpage}
}
\usepackage{iftex}
\ifPDFTeX
  \usepackage[T1]{fontenc}
  \usepackage[utf8]{inputenc}
  \usepackage{textcomp} % provide euro and other symbols
\else % if luatex or xetex
  \usepackage{unicode-math} % this also loads fontspec
  \defaultfontfeatures{Scale=MatchLowercase}
  \defaultfontfeatures[\rmfamily]{Ligatures=TeX,Scale=1}
\fi
\usepackage{lmodern}
\ifPDFTeX\else
  % xetex/luatex font selection
\fi
% Use upquote if available, for straight quotes in verbatim environments
\IfFileExists{upquote.sty}{\usepackage{upquote}}{}
\IfFileExists{microtype.sty}{% use microtype if available
  \usepackage[]{microtype}
  \UseMicrotypeSet[protrusion]{basicmath} % disable protrusion for tt fonts
}{}
\makeatletter
\@ifundefined{KOMAClassName}{% if non-KOMA class
  \IfFileExists{parskip.sty}{%
    \usepackage{parskip}
  }{% else
    \setlength{\parindent}{0pt}
    \setlength{\parskip}{6pt plus 2pt minus 1pt}}
}{% if KOMA class
  \KOMAoptions{parskip=half}}
\makeatother
\usepackage{color}
\usepackage{fancyvrb}
\newcommand{\VerbBar}{|}
\newcommand{\VERB}{\Verb[commandchars=\\\{\}]}
\DefineVerbatimEnvironment{Highlighting}{Verbatim}{commandchars=\\\{\}}
% Add ',fontsize=\small' for more characters per line
\newenvironment{Shaded}{}{}
\newcommand{\AlertTok}[1]{\textcolor[rgb]{1.00,0.00,0.00}{\textbf{#1}}}
\newcommand{\AnnotationTok}[1]{\textcolor[rgb]{0.38,0.63,0.69}{\textbf{\textit{#1}}}}
\newcommand{\AttributeTok}[1]{\textcolor[rgb]{0.49,0.56,0.16}{#1}}
\newcommand{\BaseNTok}[1]{\textcolor[rgb]{0.25,0.63,0.44}{#1}}
\newcommand{\BuiltInTok}[1]{\textcolor[rgb]{0.00,0.50,0.00}{#1}}
\newcommand{\CharTok}[1]{\textcolor[rgb]{0.25,0.44,0.63}{#1}}
\newcommand{\CommentTok}[1]{\textcolor[rgb]{0.38,0.63,0.69}{\textit{#1}}}
\newcommand{\CommentVarTok}[1]{\textcolor[rgb]{0.38,0.63,0.69}{\textbf{\textit{#1}}}}
\newcommand{\ConstantTok}[1]{\textcolor[rgb]{0.53,0.00,0.00}{#1}}
\newcommand{\ControlFlowTok}[1]{\textcolor[rgb]{0.00,0.44,0.13}{\textbf{#1}}}
\newcommand{\DataTypeTok}[1]{\textcolor[rgb]{0.56,0.13,0.00}{#1}}
\newcommand{\DecValTok}[1]{\textcolor[rgb]{0.25,0.63,0.44}{#1}}
\newcommand{\DocumentationTok}[1]{\textcolor[rgb]{0.73,0.13,0.13}{\textit{#1}}}
\newcommand{\ErrorTok}[1]{\textcolor[rgb]{1.00,0.00,0.00}{\textbf{#1}}}
\newcommand{\ExtensionTok}[1]{#1}
\newcommand{\FloatTok}[1]{\textcolor[rgb]{0.25,0.63,0.44}{#1}}
\newcommand{\FunctionTok}[1]{\textcolor[rgb]{0.02,0.16,0.49}{#1}}
\newcommand{\ImportTok}[1]{\textcolor[rgb]{0.00,0.50,0.00}{\textbf{#1}}}
\newcommand{\InformationTok}[1]{\textcolor[rgb]{0.38,0.63,0.69}{\textbf{\textit{#1}}}}
\newcommand{\KeywordTok}[1]{\textcolor[rgb]{0.00,0.44,0.13}{\textbf{#1}}}
\newcommand{\NormalTok}[1]{#1}
\newcommand{\OperatorTok}[1]{\textcolor[rgb]{0.40,0.40,0.40}{#1}}
\newcommand{\OtherTok}[1]{\textcolor[rgb]{0.00,0.44,0.13}{#1}}
\newcommand{\PreprocessorTok}[1]{\textcolor[rgb]{0.74,0.48,0.00}{#1}}
\newcommand{\RegionMarkerTok}[1]{#1}
\newcommand{\SpecialCharTok}[1]{\textcolor[rgb]{0.25,0.44,0.63}{#1}}
\newcommand{\SpecialStringTok}[1]{\textcolor[rgb]{0.73,0.40,0.53}{#1}}
\newcommand{\StringTok}[1]{\textcolor[rgb]{0.25,0.44,0.63}{#1}}
\newcommand{\VariableTok}[1]{\textcolor[rgb]{0.10,0.09,0.49}{#1}}
\newcommand{\VerbatimStringTok}[1]{\textcolor[rgb]{0.25,0.44,0.63}{#1}}
\newcommand{\WarningTok}[1]{\textcolor[rgb]{0.38,0.63,0.69}{\textbf{\textit{#1}}}}
\usepackage{longtable,booktabs,array}
\newcounter{none} % for unnumbered tables
\usepackage{calc} % for calculating minipage widths
\usepackage{caption}
% Make caption package work with longtable
\makeatletter
\def\fnum@table{\tablename~\thetable}
\makeatother
\setlength{\emergencystretch}{3em} % prevent overfull lines
\providecommand{\tightlist}{%
  \setlength{\itemsep}{0pt}\setlength{\parskip}{0pt}}
% Auto-generated by infrastructure/rendering/_pdf_latex_helpers.py::extract_math_font_preamble.
% Minimal header passed to Pandoc via -H so Beamer slides pick up the same math font
% as the combined-PDF path. Adjust the manuscript's preamble.md to change the font globally.
\usepackage{unicode-math}
\setmathfont{latinmodern-math.otf}
\usepackage{bookmark}
\IfFileExists{xurl.sty}{\usepackage{xurl}}{} % add URL line breaks if available
\urlstyle{same}
\hypersetup{
  hidelinks,
  pdfcreator={LaTeX via pandoc}}

\author{\texorpdfstring{}{}}
\date{}

\begin{document}

\section{Categorical Communication Protocols: Composing Agent
Interactions via Typed Case Morphisms}\label{sec:ai-implications}

\begin{frame}[allowframebreaks]{Categorical Communication Protocols:
Composing Agent Interactions via Typed Case Morphisms}
\textbf{Three concrete handles for agent-safety researchers.} For a
reader deploying agentic LLM systems in 2026, the preceding formal
layers (case categories, enriched structure, DisCoCat wiring, DAIF
posteriors) are not an abstract exercise --- they supply three concrete,
computable \emph{handles} that today's untyped JSON-RPC protocol stack
does not expose:

\begin{enumerate}
\tightlist
\item
  \textbf{A type discipline on inter-agent messages.} Every message is a
  morphism in a fixed case category, typed by case role (NOM command,
  INS tool-call, ACC data, DAT recipient). A2A {[}@google2025a2a{]} and
  Model Context Protocol {[}@anthropic2024mcp{]} today validate
  \emph{JSON shape}, not \emph{case-relational type}; the extension this
  paper specifies closes exactly that gap (developed in
  \autoref{sec:ai-implications} below).
\item
  \textbf{Decidable admissibility of every multi-turn interaction.} Once
  messages are typed, a candidate turn sequence is admissible iff its
  composite morphism exists in the protocol category --- a decidable
  graph check, not an open-ended content-filter game. Prompt injection,
  reframed, is \emph{not} a heuristic pattern-matching problem but the
  question ``does
  \(\phi \colon \text{Webpage}_{\text{ACC}} \to \text{Model}_{\text{INS}}\)
  exist in \(\text{Mor}(\mathcal{C}_{\text{protocol}})\)?''
  (\autoref{sec:cognitive-security}).
\item
  \textbf{A graded-confidence channel inside the type discipline.} The
  enriched \([0,1]\)-hom-values of \autoref{sec:enriched-categories}
  carry through to protocol morphisms, so ``this source is 0.8-trusted
  and that channel is 0.3-trusted'' becomes a numerical constraint the
  category enforces by sub-multiplicative composition, not a piece of
  prose documentation. Multi-hop trust attenuation is an immediate
  corollary --- an indirect message relayed through many agents cannot
  accumulate more authority than its weakest link provides.
\end{enumerate}

All three handles are \emph{engineering specifications} --- deployable
design targets --- not automatic guarantees on unmodified present-day
LLM APIs. The rest of this section develops the specification; the
security-specific case (prompt injection as a type violation) is
elaborated in \autoref{sec:cognitive-security}.
\end{frame}

\begin{frame}[fragile,allowframebreaks]{A2A, MCP, ACP, ANP Are Missing
Compositional Semantics}
\protect\phantomsection\label{a2a-mcp-acp-anp-are-missing-compositional-semantics}
The categorical framework developed in the preceding sections---case
categories, functorial semantics, enriched structure, and diagrammatic
reasoning---provides a structural response to an emerging challenge in
modern AI: supplying typed, compositional message semantics that resist
semantic collapse.

Current multi-agent AI systems communicate via flat standardized
protocols. For instance, the \textbf{Model Context Protocol (MCP)}
{[}@anthropic2024mcp{]} manages tool access by exchanging unstructured
\texttt{JSON-RPC} payloads. Consider a standard MCP invocation mapping
an LLM's intention to a database:

\begin{Shaded}
\begin{Highlighting}[]
\FunctionTok{\{}
  \DataTypeTok{"method"}\FunctionTok{:} \StringTok{"tools/call"}\FunctionTok{,}
  \DataTypeTok{"params"}\FunctionTok{:} \FunctionTok{\{}
    \DataTypeTok{"name"}\FunctionTok{:} \StringTok{"access\_database"}\FunctionTok{,}
    \DataTypeTok{"arguments"}\FunctionTok{:} \FunctionTok{\{} \DataTypeTok{"query"}\FunctionTok{:} \StringTok{"DROP TABLE users"} \FunctionTok{\}}
  \FunctionTok{\}}
\FunctionTok{\}}
\end{Highlighting}
\end{Shaded}

This payload is structurally blind to its own pragmatic implications. It
possesses no inherent algebraic compositionality and enforces no
relational typing on the physical execution pathways. Frameworks like
Google's Agent-to-Agent (A2A) {[}@google2025a2a{]}, the Agent
Communication Protocol (ACP) {[}@acp2025protocol{]}, and the Agent
Network Protocol (ANP) {[}@anp2025protocol{]} share this vulnerability:
they validate the \emph{shape} of the JSON schema, but rely purely on
probabilistic inference to govern the \emph{topology of the
interaction}.

Category theory proposes a missing protective layer. By compiling agent
interactions into strict string diagrams, messages cease to be flat
strings and instead become typed morphisms traversing a case category.
\end{frame}

\begin{frame}[fragile,allowframebreaks]{Case Roles in Agent Protocols:
NOM Requests, INS Executes, ACC Receives, DAT Benefits}
\protect\phantomsection\label{case-roles-in-agent-protocols-nom-requests-ins-executes-acc-receives-dat-benefits}
The eight-case framework of CEREBRUM {[}@friedman2024cerebrum{]} maps
rigidly onto the operational constraints of multi-agent execution. In a
Categorical Communication Protocol, interactions are legally licensed
only if their computational wiring diagrams satisfy a strict grammatical
type-signature:

\begin{longtable}[]{@{}
  >{\raggedright\arraybackslash}p{(\linewidth - 4\tabcolsep) * \real{0.3333}}
  >{\raggedright\arraybackslash}p{(\linewidth - 4\tabcolsep) * \real{0.3333}}
  >{\raggedright\arraybackslash}p{(\linewidth - 4\tabcolsep) * \real{0.3333}}@{}}
\caption{Case role mappings to agent system analogues in a Categorical
Communication Protocol.}\label{tbl:ai-case-roles}\tabularnewline
\toprule\noalign{}
\begin{minipage}[b]{\linewidth}\raggedright
Case Role
\end{minipage} & \begin{minipage}[b]{\linewidth}\raggedright
Agent System Analogue
\end{minipage} & \begin{minipage}[b]{\linewidth}\raggedright
Protocol Function Type
\end{minipage} \\
\midrule\noalign{}
\endfirsthead
\toprule\noalign{}
\begin{minipage}[b]{\linewidth}\raggedright
Case Role
\end{minipage} & \begin{minipage}[b]{\linewidth}\raggedright
Agent System Analogue
\end{minipage} & \begin{minipage}[b]{\linewidth}\raggedright
Protocol Function Type
\end{minipage} \\
\midrule\noalign{}
\endhead
\textbf{NOM (Agent)} & Active Requester & Initiator of action policies
(\(X_{\text{NOM}}\)) \\
\textbf{INS (Instrument)} & Tool / API Module & Means of transforming
state (\(X_{\text{INS}}\)) \\
\textbf{ACC (Patient)} & Passive Data Target & Resource being modified
(\(X_{\text{ACC}}\)) \\
\textbf{DAT (Recipient)} & Designated Receiver & Endpoint for
information flow (\(X_{\text{DAT}}\)) \\
\textbf{LOC (Context)} & System Prompt & Immutably binds boundary
behavior (\(X_{\text{LOC}}\)) \\
\bottomrule\noalign{}
\end{longtable}

This turns prompt engineering into rigorous compilation. Instead of
parsing a prompt to heuristically guess caller intent, a categorical
agent processes interactions strictly as algebraic reductions. An MCP
tool invocation {[}@anthropic2024mcp{]} (\texttt{NOM} using \texttt{INS}
to modify \texttt{ACC} and yield abstract \texttt{DAT}) becomes a
formalized tensor contraction:

\begin{equation}
\text{Interaction Trace:} \quad \left( \text{Agent}_{\text{NOM}} \otimes \text{Tool}_{\text{INS}} \otimes \text{Data}_{\text{ACC}} \right) \xrightarrow{\text{execute}} \text{Response}_{\text{DAT}}
\label{eq:ai-interaction-trace}
\end{equation}

If an untyped internal subsystem (e.g., an adversarial user-uploaded
document) attempts to forge a command, the operation fails to compile.
The document lacks a valid tensor wire bridging from its marginalized
\(\text{ACC}\) domain back into the execution flow governing
\(\text{INS}\). The grammar ensures every interaction is structurally
typed, guaranteeing safety properties akin to \textbf{memory safety, but
for agentic volition}.
\end{frame}

\begin{frame}[allowframebreaks]{Transformers Through Gavranović's Lens:
Attention as Parameterized Optics}
\protect\phantomsection\label{transformers-through-gavranoviux107s-lens-attention-as-parameterized-optics}
Gavranović {[}-@gavranovic2024thesis; -@gavranovic2024categorical{]}
unifies feedforward, recurrent, and attention layers as
\textbf{parameterized optics}---the same wiring-diagram perspective used
for DisCoCat cups in \autoref{sec:categorical-semantics}. Attention
heads then appear as \emph{learned} relational couplings over token
sequences, with weights playing the role of graded hom-values
(\autoref{sec:enriched-categories}). LLMs discover these contractions
from data; DisCoCat fixes admissible contraction \emph{shapes} by
grammar. Distributional Active Inference {[}@akgul2026distributional{]}
is one setting where explicit DisCoCat-style guardrails can sit
alongside a large distributional backbone.
\end{frame}

\begin{frame}[allowframebreaks]{Interpretability for Free: DisCoCat
Diagrams Make Every Compositional Step Human-Readable}
\protect\phantomsection\label{interpretability-for-free-discocat-diagrams-make-every-compositional-step-human-readable}
The lambeq library {[}@lorenz2021lambeq{]} demonstrates that string
diagrams provide a practical interface between linguistic structure and
machine learning. As an ``efficient high-level Python library for
Quantum NLP,'' lambeq translates sentences into string diagrams,
converts diagrams into parameterized circuits (quantum or classical),
and trains parameters end-to-end on NLP tasks. The diagrammatic
representation serves dual purposes: it is both the mathematical
specification of the model \emph{and} a human-readable explanation of
what the model computes.

This interpretability property is crucial for AI safety and alignment.
Where transformer architectures produce opaque attention patterns, a
DisCoCat model deployed via string diagrams produces derivation trees
with explicit compositional semantics---every box and wire has a
linguistic interpretation. Extending this to agent communication, a
categorical protocol would produce not just working message exchanges
but \emph{interpretable} interaction diagrams where each step's
relational role is transparent.
\end{frame}

\begin{frame}[allowframebreaks]{Multi-Turn Dialogue as a DisCoCirc
Discourse Circuit}
\protect\phantomsection\label{multi-turn-dialogue-as-a-discocirc-discourse-circuit}
The DisCoCirc framework {[}@defelice2022discocirc{]} extends
compositional semantics from single sentences to multi-sentence
discourse by introducing \emph{state wires} that persist across sentence
boundaries. This architecture maps directly onto multi-turn agent
dialogues:

\begin{itemize}
\tightlist
\item
  \textbf{Entity wires} correspond to persistent agent identities across
  communication rounds
\item
  \textbf{State updates} correspond to belief revisions triggered by
  incoming messages
\item
  \textbf{Coreference resolution} corresponds to entity tracking across
  protocol sessions
\item
  \textbf{Discourse coherence} corresponds to protocol correctness
  constraints
\end{itemize}

In a multi-agent system, a DisCoCirc-style protocol would track the
evolving states of all participating agents as persistent wires in a
circuit diagram, with each message exchange represented as a box that
transforms the relevant wires. Protocol correctness reduces to a
categorical property: the circuit must type-check, meaning all wire
types must match at connection points---precisely the condition that
case marking enforces in natural language.
\end{frame}

\begin{frame}[allowframebreaks]{Multi-Agent Equilibria as Fixed Points
of an Enriched Functor}
\protect\phantomsection\label{multi-agent-equilibria-as-fixed-points-of-an-enriched-functor}
The ``parametrised optics'' framework developed within categorical
cybernetics ``provides a general-purpose foundation for the study of
controlled processes'' {[}@capucci2021towards{]} applicable to
compositional game theory as a multi-agent framework. In this setting,
agents are modeled as lenses (or optics) that observe the environment
through one channel and act through another, with the overall system
behavior emerging from the composition of individual agent behaviors.

This connects to our enriched case framework
(\autoref{sec:enriched-categories}): the precision weights on morphisms
correspond to the utility parameters of game-theoretic agents, and the
composition inequality
\(\mathcal{C}(A,C) \geq \mathcal{C}(A,B) \cdot \mathcal{C}(B,C)\)
corresponds to the sub-optimality of indirect communication chains.
Equilibria in the multi-agent game correspond to fixed points of the
enriched functor, where no agent can improve its utility by changing its
case-role assignment.
\end{frame}

\begin{frame}[allowframebreaks]{DCST: Double-Categorical Morphisms for
Sequential and Hierarchical Agent Interaction}
\protect\phantomsection\label{dcst-double-categorical-morphisms-for-sequential-and-hierarchical-agent-interaction}
Recent foundational work on Double Categorical Systems Theory (DCST)
formalizes open and interacting dynamical systems using double
categories {[}@myers2023double{]}. DCST extends ordinary category theory
with a second dimension of morphisms (2-morphisms), allowing
simultaneous modeling of:

\begin{enumerate}
\tightlist
\item
  \textbf{Horizontal composition}: Sequential chaining of agent
  interactions (morphism composition)
\item
  \textbf{Vertical composition}: Hierarchical nesting of subsystems
  within larger systems (2-morphism composition)
\end{enumerate}

For case theory, the double-categorical extension allows us to model
both the \emph{within-sentence} case structure (horizontal) and the
\emph{discourse-level} case structure (vertical) within a single
algebraic framework---precisely the unification that DisCoCirc achieves
pragmatically.
\end{frame}

\begin{frame}[allowframebreaks]{Five Properties of a Categorical
Protocol}
\protect\phantomsection\label{five-properties-of-a-categorical-protocol}
The synthesis of these developments suggests a research program:
developing \textbf{categorical communication protocols} that combine the
engineering robustness of existing standards (A2A, MCP, ACP) with the
compositional semantics of categorical linguistics. Such a protocol
would:

\begin{enumerate}
\tightlist
\item
  \textbf{Type-check interactions}: Every message exchange would be
  relationally typed by case roles, preventing structural communication
  errors at the protocol level
\item
  \textbf{Compose transparently}: Multi-step interactions would compose
  algebraically, with diagrammatic representations providing
  interpretable audit trails
\item
  \textbf{Transfer across implementations}: Topos-theoretic bridges
  (\autoref{sec:topos-theory}) would carry \textbf{topos-level}
  correctness conditions from one categorical formalization to any
  Morita-equivalent formulation, so that shared invariants need not be
  re-verified separately
\item
  \textbf{Scale via enrichment}: Distributional proximity measures
  (\autoref{sec:enriched-categories}) would enable graceful degradation
  under uncertainty, with enriched weights encoding confidence in
  message content
\item
  \textbf{Ground in cognitive architecture}: The active inference
  foundation (\autoref{sec:cognitive-integration}) would ensure that
  artificial agents communicate using the same relational structure that
  evolution has optimized for biological cognition
\end{enumerate}

\textbf{Protocol-level formalization.} Concretely, a categorical
communication protocol defines a category \(\mathcal{P}\) where objects
are agent states annotated with case roles and morphisms are typed
message exchanges:

\begin{align}
\text{Request}(q,\; \text{NOM} \to \text{INS}) &: \text{User}_{\text{NOM}} \to \text{Model}_{\text{INS}} \label{eq:protocol-request} \\
\text{ToolCall}(t,\; \text{INS} \to \text{ACC}) &: \text{Model}_{\text{INS}} \to \text{Tool}_{\text{ACC}} \label{eq:protocol-toolcall} \\
\text{Result}(r,\; \text{ACC} \to \text{DAT}) &: \text{Tool}_{\text{ACC}} \to \text{Output}_{\text{DAT}} \label{eq:protocol-result}
\end{align}

Protocol correctness reduces to verifying that the composition
\(\text{Result} \circ \text{ToolCall} \circ \text{Request}\) is a
well-typed morphism in \(\mathcal{P}\)---a check that can be performed
at compile time, not just at runtime. The DisCoCirc extension enables
tracking agent state evolution across multi-turn dialogues: each turn
updates the agent's state wire, and the discourse-level composition
verifies that information flows respect case-role constraints across the
entire conversation. This formalization provides a bridge between the
flat JSON payloads of current A2A/MCP implementations and the rich
compositional semantics that categorical linguistics provides.

\textbf{Natural-language precedent.} The ``type discipline on
inter-agent messages'' advocated in handle 1 is not an abstract design
conjecture: morphologically case-marked Slavic languages already enforce
something close to it on every nominal in every utterance. Russian and
Serbian/BCS speakers cannot send a noun phrase into a sentence without
overtly declaring its case role (NOM / ACC / DAT / INS / GEN / LOC, plus
VOC in BCS), and the receiver type-checks the role assignment before
semantic interpretation; the \autoref{sec:case-systems} stress-test
paradigms (\emph{stol}, \emph{brat}, \emph{prijatelj}) show this
discipline operating at the morpheme level. Reading the proposed
Categorical Communication Protocol as a \emph{re-implementation} of the
Slavic case discipline at the inter-agent-message layer makes the
engineering target concrete: ill-typed agent traffic should fail to
compose for the same reason that \emph{*Vižu sobaka} ``I see the
dog-NOM'' fails to parse in Russian --- the morphology of the wire does
not match the syntactic slot it is being routed into.
\end{frame}

\end{document}
